\section{Absolute Stieltjes reconstruction and logarithmic dilation identities}
\label{stconf:section}

Repeated zeros in the dilation polynomial control logarithmic moments
as well as ordinary power moments. At the zero-Mellin endpoint this
gives a direct reconstruction of the \emph{absolute} generalized
Stieltjes functions from convergent digamma integrals. The translation
formulas in Section~\ref{stmell:section} reconstruct their differences;
the additional dilation cancellation below removes both endpoint
obstructions without a split-integral subtraction. The scalar Mellin
foundation remains the one proved in \cite{stmell:balanced} and
recalled in Lemma~\ref{stbil:scalar}. The higher-logarithm coefficient
principle is already present in that report; the results here specify
the dilation kernels, their endpoint convergence, and an invertible
coefficient formula.

\subsection{An ordinary digamma integral for every absolute Stieltjes index}

Let $a>0$ and $L,M>0$, and set $p=LM$. Define
\begin{align}
 \mathcal K_{L,M}(x,a)
   ={}&-\psi\left(a+e^{(L+M)/2}x\right)
       +\psi\left(a+e^{(M-L)/2}x\right)\notag\\
     &+\psi\left(a+e^{(L-M)/2}x\right)
       -\psi\left(a+e^{-(L+M)/2}x\right),\qquad x>0,
       \label{stconf:digamma-kernel}\\
 N_j(a;L,M)&=\int_0^\infty
           \frac{\log^j x}{x}\mathcal K_{L,M}(x,a)\,dx.
       \label{stconf:N}
\end{align}
Thus equal scales give the three-term expression
\[
 \mathcal K_{L,L}(x,a)
      =2\psi(a+x)-\psi(a+e^Lx)-\psi(a+e^{-L}x).
\]
Introduce the even holomorphic germ
\begin{equation}\label{stconf:even-germ}
 E_{L,M}(u)=\frac{\sin\pi u}{\pi u}
       \frac{Lu}{2\sinh(Lu/2)}\frac{Mu}{2\sinh(Mu/2)}
       =\sum_{j\ge0}e_{2j}(L,M)u^{2j}.
\end{equation}
Its constant term is one. All appearances of these quotients at zero
mean their removable values.

\begin{theorem}[Absolute digamma reconstruction at every index]
\label{stconf:absolute}
Every integral \eqref{stconf:N} converges absolutely. For all $n\ge0$,
\begin{equation}\label{stconf:absolute-formula}
 \boxed{\displaystyle
 \gamma_n(a)=\frac{n!}{LM}
   \sum_{j=0}^{\lfloor(n+1)/2\rfloor}
    \frac{e_{2j}(L,M)}{(n+1-2j)!}N_{n+1-2j}(a;L,M).}
\end{equation}
In particular $N_0=-LM$. Equivalently, with the explicit polynomial
\[
 P_n(y;L,M)=n!\sum_{j=0}^{\lfloor(n+1)/2\rfloor}
                    \frac{e_{2j}(L,M)y^{n+1-2j}}{(n+1-2j)!},
\]
each coefficient is a single ordinary integral:
\begin{equation}\label{stconf:single}
 \gamma_n(a)=\frac1{LM}\int_0^\infty
          P_n(\log x;L,M)\mathcal K_{L,M}(x,a)\,\frac{dx}{x}.
\end{equation}
No choice of an integration cutoff or finite-part constant is involved.
\end{theorem}
\begin{proof}
At zero the four constant terms $\psi(a)$ cancel, so the kernel is
$O(x)$. At infinity use
$\psi(a+bx)=\log x+\log b+O(x^{-1})$ for each positive $b$.
Both logarithmic combinations cancel, so the kernel is $O(x^{-1})$.
Consequently its Mellin integral is holomorphic on
$-1<\Re u<1$. On every smaller closed strip, each derivative in the Mellin parameter
is dominated by the same endpoint bounds times a fixed power of
$1+|\log x|$. This proves ordinary absolute convergence of every
$N_j$ and permits differentiation under the integral.

The relation $C_0(x,a)=\log x-\psi(x+a)$ writes the kernel as
\[
 C_0(e^{(L+M)/2}x,a)-C_0(e^{(M-L)/2}x,a)
 -C_0(e^{(L-M)/2}x,a)+C_0(e^{-(L+M)/2}x,a).
\]
On $0<\Re u<1$ its scalar Mellin transform is therefore
\begin{equation}\label{stconf:digamma-mellin}
 \int_0^\infty x^{u-1}\mathcal K_{L,M}(x,a)\,dx
 =4\sinh(Lu/2)\sinh(Mu/2)
                \frac{\pi}{\sin\pi u}\zeta(1-u,a).
\end{equation}
The independently proved endpoint bounds continue this identity
holomorphically through $u=0$. Multiplication by
\eqref{stconf:even-germ} gives the local equality
\begin{equation}\label{stconf:absolute-germ}
 \frac{E_{L,M}(u)}{LM}\sum_{j\ge0}N_j\frac{u^j}{j!}
       =u\zeta(1-u,a)
       =-1+\sum_{n\ge0}\gamma_n(a)\frac{u^{n+1}}{n!}.
\end{equation}
The series represent holomorphic germs near zero, rather than merely
formal expressions for divergent integrals. Constant and higher
coefficients prove all the claims.
\end{proof}

The coefficients in \eqref{stconf:even-germ} have the finite formula
\begin{equation}\label{stconf:coefficient}
 e_{2j}(L,M)=
 \sum_{v+w+t=j}
 \frac{(-1)^v\pi^{2v}}{(2v+1)!}
 \frac{B_{2w}(1/2)L^{2w}}{(2w)!}
 \frac{B_{2t}(1/2)M^{2t}}{(2t)!},
\end{equation}
by the classical Bernoulli generating function. In particular put
\[
 A=\frac{\pi^2}{6}+\frac{L^2+M^2}{24};
 \qquad e_0=1,\quad e_2=-A.
\]
The first three absolute reconstruction formulas are
\begin{align}
 \gamma_0(a)&=\frac{N_1}{LM},\label{stconf:gamma0}\\
 \gamma_1(a)&=\frac{N_2}{2LM}+A,\label{stconf:gamma1}\\
 \gamma_2(a)&=\frac{N_3}{3LM}-2A\frac{N_1}{LM}.
                   \label{stconf:gamma2}
\end{align}
These use the conventional Laurent coefficients, so
$\gamma_0(a)=-\psi(a)$. Equal scales give
$A=\pi^2/6+L^2/12$. Odd-index coefficients use even logarithmic
moments and the known constant $N_0$; even-index coefficients use
only odd logarithmic moments.

For comparison with uncentered dilations, put $\lambda=e^L$,
$\mu=e^M$ and
\[
 H(x,a)=-\psi(a+x)+\psi(a+x/\lambda)+\psi(a+x/\mu)
                             -\psi(a+x/(\lambda\mu)).
\]
Then $\mathcal K_{L,M}(x,a)=H(\sqrt{\lambda\mu}\,x,a)$.
Thus $N_j$ is exactly the ordinary integral of
$H(x,a)(\log x-\tfrac12\log(\lambda\mu))^j/x$.
The symmetry and the even coefficient germ arise from this change
of variable; no endpoint term is discarded.

\subsection{The colored endpoint and an alternating example}

Let $\omega\in\mathbb C\setminus[1,\infty)$ and $a,L>0$. With the
Lerch branch and colored kernels of Section~\ref{stcol:section}, set
\begin{align}
 \mathcal J_L^\omega(x,a)&=
      \Phi(\omega,1,a+e^{L/2}x)
          -\Phi(\omega,1,a+e^{-L/2}x),\notag\\
 T_j&=\int_0^\infty\frac{\log^j x}{x}
                                 \mathcal J_L^\omega(x,a)\,dx,\notag\\
 E_L^{\rm col}(u)&=\frac{\sin\pi u}{\pi u}
                         \frac{Lu}{2\sinh(Lu/2)}
                    =\sum_{j\ge0}f_{2j}(L)u^{2j}.
\end{align}
\begin{theorem}[Colored absolute endpoint reconstruction]
\label{stconf:colored-endpoint}
Every $T_j$ converges absolutely, and for every $n\ge0$,
\begin{equation}\label{stconf:colored-inversion}
 \partial_s^n\Phi(\omega,1,a)
 =\frac{(-1)^{n+1}n!}{L}
     \sum_{j=0}^{\lfloor n/2\rfloor}
             \frac{f_{2j}(L)}{(n-2j)!}T_{n-2j}.
\end{equation}
\end{theorem}
\begin{proof}
Analyticity at the positive shift gives $\mathcal J=O(x)$ at zero.
The colored tail expansion gives $\mathcal J=O(x^{-1})$ at infinity.
The same domination argument as above gives a holomorphic Mellin
integral on $-1<\Re u<1$. Scaling in the colored scalar transform
first on $0<\Re u<1$, then continuing through zero, gives
\[
 \sum_{j\ge0}T_j\frac{u^j}{j!}
  =-2\sinh(Lu/2)\frac\pi{\sin\pi u}\Phi(\omega,1-u,a).
\]
Multiplying by $E_L^{\rm col}$ gives $-L\Phi(\omega,1-u,a)$.
Coefficient extraction proves \eqref{stconf:colored-inversion}.
The absence of the uncolored Laurent pole explains why one dilation
difference suffices here, whereas the digamma reconstruction uses two.
\end{proof}

For $\omega=-1$, $a=1$, and $L=\log2$, put
\[
 f(x)=\tfrac12\left[\psi\left(\tfrac{x+2}{2}\right)
                        -\psi\left(\tfrac{x+1}{2}\right)\right],
 \qquad J(x)=f(\sqrt2x)-f(x/\sqrt2).
\]
The identity $\Phi(-1,s,1)=(1-2^{1-s})\zeta(s)$ and
$f_2=-\pi^2/6-L^2/24$ give the explicit ordinary moments
\begin{align}
 \int_0^\infty J(x)\,\frac{dx}{x}&=-L^2,\notag\\
 \int_0^\infty \log x\,J(x)\,\frac{dx}{x}
     &=L^2\gamma-\tfrac12L^3,\notag\\
 \int_0^\infty \log^2x\,J(x)\,\frac{dx}{x}
     &=2L^2\gamma_1+L^3\gamma-\frac{\pi^2L^2}{3}
                                       -\frac{5L^4}{12}.
        \label{stconf:alternating-example}
\end{align}
For any root-of-unity color the same theorem is a generalized-Stieltjes
identity through the finite dictionary \eqref{stcol:dictionary}.

\subsection{Repeated roots at ordinary positive-integer moments}

The same finite-dilation mechanism applies either to the uncolored
kernel $\mathcal C_q=C_q$ with $G(s)=\zeta(s,a)$, or the colored
kernel $\mathcal C_q=C_q^\omega$ with $G(s)=\Phi(\omega,s,a)$.
The positive shifts and color domain are as above. Fix integers
$0\le q\le K$. For each $k=q,\ldots,K$ choose a multiplicity
$r_k\ge1$ and numbers $\lambda_{k,\nu}>1$, $1\le\nu\le r_k$.
Put
\begin{equation}\label{stconf:product}
 D(s)=\prod_{k=q}^K\prod_{\nu=1}^{r_k}
                  (1-\lambda_{k,\nu}^{-s-k})
        =\sum_h c_hh^{-s},\qquad
 \mathcal H_q(x)=\sum_h\frac{c_h}{h}\mathcal C_q(x/h,a).
\end{equation}
Repeated scales in the finite expansion are combined. Its Laurent
subtraction in the uncolored case remains $D(1)/z$.

\begin{theorem}[Confluent ordinary moment blocks]
\label{stconf:blocks}
For every $k=q,\ldots,K$ and $\ell\ge0$, the integral
\[
 M_{k,\ell}=\int_0^\infty x^k\log^\ell x\,\mathcal H_q(x)\,dx
\]
is absolutely convergent. Let $\rho_k$ be the zero multiplicity of
$G$ at $-k$, with $\rho_k=0$ when $G(-k)\ne0$. Then
\begin{equation}\label{stconf:vanishing-order}
 M_{k,\ell}=0\quad(0\le\ell<r_k+\rho_k-1),\qquad
 M_{k,r_k+\rho_k-1}\ne0.
\end{equation}
In particular, put
\[
 W_k=\prod_{\nu=1}^{r_k}\log\lambda_{k,\nu}\;
      \prod_{\substack{j=q\\j\ne k}}^K
            \prod_{\nu=1}^{r_j}(1-\lambda_{j,\nu}^{k-j}).
\]
Then the first potentially nonzero logarithmic moment is
\begin{equation}\label{stconf:first-log}
 M_{k,r_k-1}
   =(-1)^{k+r_k+1}(r_k-1)!Q_q(k+1)W_kG(-k).
\end{equation}
More generally the complete convergent moment germ is
\begin{equation}\label{stconf:bulk-germ}
 \sum_{\ell\ge0}M_{k,\ell}\frac{u^\ell}{\ell!}
  =(-1)^{k+1}\frac\pi{\sin\pi u}
       Q_q(k+1+u)D(-k-u)G(-k-u).
\end{equation}
\end{theorem}
\begin{proof}
For either scalar family the tail coefficient indexed by $j\ge q$
is a nonzero Pochhammer factor times $D(-j)G(-j)$. The specified
roots cancel every index $q,\ldots,K$, so
$\mathcal H_q(x)=O(x^{-K-2})$. At zero it is bounded by a fixed
power of $1+|\log x|$. Consequently all the displayed moments are
ordinary and absolutely convergent. For $u$ in a sufficiently small
disc, their integrands and all derivatives are dominated at infinity
by $x^{-2+\epsilon}$ times powers of $\log x$, and at zero by
$x^{k-\epsilon}$ times such powers. This proves the holomorphic
moment germ and its coefficient interpretation.

Finite summation of the appropriate scalar Mellin transform identifies
that germ with \eqref{stconf:bulk-germ}. At $-k$, the factors assigned
to $k$ have simple zeros, and none assigned to another integer vanishes.
Thus $D(-k-u)=(-1)^{r_k}W_ku^{r_k}+O(u^{r_k+1})$.
Also $Q_q(k+1)=q!e_q(1,1/2,\ldots,1/k)>0$ since $k\ge q$.
The sine contributes exactly one pole. The total order is therefore
$r_k+\rho_k-1$, proving the sharp vanishing assertion, while its
leading coefficient for $\rho_k=0$ gives \eqref{stconf:first-log}.
The scalar functions here are not identically zero, so their zero
orders are finite.
\end{proof}

For a single repeated scale at $k=q$, take
\[
 \mathcal H_q(x)=\sum_{j=0}^r(-1)^j\binom rj
                 \lambda^{-(q+1)j}\mathcal C_q(x/\lambda^j,a).
\]
Its first $r-1$ logarithmic moments of weight $x^q$ vanish and
\[
 \int_0^\infty x^q\log^{r-1}x\,\mathcal H_q(x)\,dx
   =(-1)^{q+r+1}(r-1)!(\log\lambda)^rG(-q).
\]
For example, at the alternating color and $q=1,r=2$,
\[
 V(x)=C_1^{-1}(x,a)-2\lambda^{-2}C_1^{-1}(x/\lambda,a)
                              +\lambda^{-4}C_1^{-1}(x/\lambda^2,a)
\]
satisfies the exact pair
\begin{equation}\label{stconf:colored-bulk}
 \int_0^\infty xV(x)\,dx=0,\qquad
 \int_0^\infty x\log x\,V(x)\,dx
                     =\frac{2a-1}{4}(\log\lambda)^2.
\end{equation}
Here $\Phi(-1,-1,a)=(2a-1)/4$; at $a=1/2$ the second moment
also vanishes, exactly as the sharper zero-order statement predicts.

\subsection{The support bound for prescribed multiplicities}

Let $R=\sum_{k=q}^Kr_k$. Any nonzero Dirichlet polynomial
$D(s)=\sum_{j=1}^Jc_jh_j^{-s}$ with distinct positive scales and
zeros of multiplicity at least $r_k$ at all $-k$ requires
$J\ge R+1$. To prove this, a nonzero real exponential sum with $J$
distinct real exponents has at most $J-1$ real zeros counted with
multiplicity: divide by one exponential, differentiate, and apply
Rolle's theorem inductively. For complex coefficients, choose a nonzero
real or imaginary component; all prescribed zeros are retained in
that component. This proves the same lower bound. A common scale
$\lambda>1$ attains it, since
\[
 D(s)=\prod_{k=q}^K(1-\lambda^{-s-k})^{r_k}
\]
is a degree-$R$ polynomial in $\lambda^{-s}$ with all $R+1$
coefficients nonzero. Up to normalization it is the unique filter on
the prescribed scales $1,\lambda,\ldots,\lambda^R$ with these zero
multiplicities. This is a statement about the specified zeros of $D$;
intrinsic zeros of $G$ can supply additional moment vanishing without
increasing the number of scales.

\subsection{Principal characters and universal centered logarithmic moments}
\label{stlog:principal}

Repeated Euler zeros give particularly explicit logarithmic identities.
Let $d>1$, let $p_1,\ldots,p_r$ be its distinct prime divisors, and put
\[
 d_0=\prod_{j=1}^rp_j,\quad \ell_j=\log p_j,\quad
 A_d=\frac{(-1)^r}{2}\prod_{j=1}^r\ell_j,\quad
 c_d=\frac12\log d_0-\log(2\pi).
\]
Use the principal kernel of \eqref{stdil:principal-kernel} at order zero:
\[
 K_d(x)=\sum_{h\mid d_0}\frac{\mu(h)}h
             [\log(x/h)-\psi(1+x/h)].
\]
For $d>1$, it is $O(1+|\log x|)$ at zero and $O(x^{-2})$ at infinity.
Thus every centered logarithmic moment
\[
 N_j(d)=\int_0^\infty (\log x-c_d)^jK_d(x)\,dx
 \qquad(j\ge0)
\]
is an absolutely convergent ordinary integral.

\begin{theorem}[The principal logarithmic germ]
\label{stlog:principal-germ}
Define $\operatorname{sinhc}z=\sinh(z)/z$, with value one at zero.
In a neighborhood of $h=0$,
\begin{align}
 \sum_{j\ge0}\frac{N_j(d)}{j!}h^j
 &=A_dh^{r-1}\,\Gamma(1+h)[h\zeta(1+h)]
       \sec\frac{\pi h}{2}
       \prod_{j=1}^r\operatorname{sinhc}\frac{h\ell_j}{2}.
       \label{stlog:principal-series}
\end{align}
The apparently singular factor $h\zeta(1+h)$ has its removable value one.
Writing $\gamma=\gamma_0(1)$ and $\gamma_j=\gamma_j(1)$, the first moments are
\begin{align}
 N_j(d)&=0 &&(0\le j<r-1),\notag\\
 N_{r-1}(d)&=A_d(r-1)!,&
 N_r(d)&=0,\label{stlog:first-moments}\\
 N_{r+1}(d)&=A_d(r+1)!
  \left(\frac{5\pi^2}{24}-\frac{\gamma^2}{2}-\gamma_1
                 +\frac1{24}\sum_{j=1}^r\ell_j^2\right),\notag\\
 N_{r+2}(d)&=A_d(r+2)!
  \left(\frac{\gamma_2}{2}+\gamma\gamma_1
                          +\frac{\gamma^3-\zeta(3)}3\right).
                          \label{stlog:universal-cubic}
\end{align}
In particular, the normalized coefficient in the last line is independent
of the chosen primes and of their number.
\end{theorem}
\begin{proof}
The endpoint estimates give a holomorphic Mellin integral on
$-1<\Re h<1$, and allow differentiation of every fixed order near zero.
By \eqref{stdil:mellin}, its uncentered germ is
\[
 \int_0^\infty x^hK_d(x)\,dx
   =\pi\csc\bigl(\pi(1+h)\bigr)\zeta(-h)
                   \prod_{j=1}^r(1-p_j^h).
\]
The classical zeta reflection formula \cite{stchar:classical} gives
\[
 \zeta(-h)=-(2\pi)^{-h}\frac{\sin(\pi h/2)}{\pi}
                     \Gamma(1+h)\zeta(1+h),
\]
while the elementary product identity is
\[
 \prod_{j=1}^r(1-p_j^h)
   =(-1)^rh^r\left(\prod_j\ell_j\right)
       e^{h\sum_j\ell_j/2}
       \prod_j\operatorname{sinhc}\frac{h\ell_j}{2}.
\]
Combining them and multiplying by $e^{-c_dh}$ proves
\eqref{stlog:principal-series}, first away from zero and then by removability.
The familiar gamma Taylor expansion and the defining Stieltjes expansion yield
\begin{align*}
 \Gamma(1+h)[h\zeta(1+h)]
   &=1+\left(\frac{\zeta(2)}2-\frac{\gamma^2}2-\gamma_1\right)h^2\\
   &\quad+\left(\frac{\gamma_2}2+\gamma\gamma_1
                  +\frac{\gamma^3-\zeta(3)}3\right)h^3+O(h^4).
\end{align*}
There is no linear term. The secant and hyperbolic factors are even,
with quadratic coefficient $\pi^2/8+\sum_j\ell_j^2/24$.
Coefficient extraction proves every displayed formula. In particular,
the absence of a linear term proves the centered zero at order $r$,
and evenness leaves the cubic coefficient independent of the primes.
\end{proof}

For instance, $d=6$ gives
\[
 K_6(x)=C_0(x,1)-\tfrac12C_0(x/2,1)
                    -\tfrac13C_0(x/3,1)+\tfrac16C_0(x/6,1),
\]
and the concrete identities
\[
 \int_0^\infty K_6(x)\,dx=0,\qquad
 \int_0^\infty\log x\,K_6(x)\,dx=\frac{\log2\log3}{2},
\]
\[
 \int_0^\infty
       \log^2\!\left(\frac{2\pi x}{\sqrt6}\right)K_6(x)\,dx=0.
\]
Higher prime powers do not change these kernels or identities: only the
distinct prime divisors enter. Their source is the exact Euler product
and classical functional equation, not a conjectural arithmetic relation.
